%---------------------------------------------------------------------------%
%->> 封面信息及生成
%---------------------------------------------------------------------------%
\confidential{}
\schoollogo{scale=0.2}{ShanghaiTechLogo}
\title{基于物理信息神经网络与大语言模型的统一期权定价框架}
\author{朱文龙}
\ID{2022XXXXXXXX}
\entranceYear{2022}
\advisor{指导教师姓名}
\advisorsec{}
\degree{学士}
\degreetype{工学}
\major{计算机科学与技术}
\institute{上海科技大学信息科学与技术学院}
\chinesedate{二零二六年~六月}
\englishtitle{A Unified Option Pricing Framework Based on Physics-Informed Neural Networks and Large Language Models}
\englishauthor{Wenlong Zhu}
\englishadvisor{Supervisor: }
\englishdegree{Bachelor}
\englishdegreetype{Engineering}
\englishthesistype{thesis}
\englishmajor{Computer Science and Technology}
\englishinstitute{School of Information Science and Technology}
\englishdate{June / 2026}

%---------------------------------------------------------------------------%
%->> 中文摘要
%---------------------------------------------------------------------------%
\chapter*{摘\quad 要}\chaptermark{摘\quad 要}
\setcounter{page}{1}
\pagenumbering{Roman}

期权定价是金融工程的核心问题之一。Black-Scholes-Merton（BSM）模型、常弹性方差（CEV）模型和Heston随机波动率模型是三大经典期权定价框架，分别适用于不同的市场假设。传统数值方法（有限差分、有限元）在求解高维偏微分方程（PDE）时面临维数灾难，而现有物理信息神经网络（PINN）研究大多针对单一模型，缺乏统一框架。

本文提出并实现了一个基于PINN与大语言模型（LLM）的统一期权定价框架。在PINN部分，通过软掩码（Soft Mask）机制和加法参数化输出，将BSM、CEV、Heston三大模型的PDE算子统一为单一参数化形式，由一个9维输入的深层全连接网络同时覆盖三种模型的解映射，无需分别训练独立网络。软掩码通过$\tanh(\xi/0.05)^2$实现模型间的连续插值，当随机波动率参数$\xi\to 0$时自动退化为BSM/CEV，当$\xi$较大时激活Heston的随机波动率项。加法参数化以Black-Scholes解析解为基线，网络学习绝对修正量，解决了乘法参数化在深度虚值区域退化的问题。对Heston模型采用改进约束PINN（ICPINN）硬约束技术，通过输出变换精确满足终值条件。

在LLM部分，引入DeepSeek API作为智能路由层，通过结构化JSON输出实现可靠的参数提取和模型选择，支持中英文自然语言输入，使非专业用户无需了解模型细节即可完成期权定价。

实验结果表明，统一PINN在BSM模型上达到MAE=0.067（ATM区域相对误差0.31\%），在Heston模型上达到MAE=0.031（ATM区域相对误差0.52\%），均优于或持平于独立训练的PINN基线。特别地，独立训练的Heston PINN完全无法收敛，而统一框架成功解决了这一问题。LLM路由层在20个测试用例上实现了100\%的模型选择准确率。单次推断耗时约2ms，远低于传统有限差分法的约500ms。

\keywords{物理信息神经网络；期权定价；Black-Scholes-Merton模型；Heston模型；大语言模型；统一框架}

%---------------------------------------------------------------------------%
%->> 英文摘要
%---------------------------------------------------------------------------%
\chapter*{Abstract}\chaptermark{Abstract}

Option pricing is one of the core problems in financial engineering. The Black-Scholes-Merton (BSM) model, the Constant Elasticity of Variance (CEV) model, and the Heston stochastic volatility model are three classical option pricing frameworks applicable to different market assumptions. Traditional numerical methods (finite difference, finite element) suffer from the curse of dimensionality when solving high-dimensional partial differential equations (PDEs), while existing Physics-Informed Neural Network (PINN) research mostly targets single models and lacks a unified framework.

This paper proposes and implements a unified option pricing framework based on PINN and Large Language Models (LLM). In the PINN component, a Soft Mask mechanism and additive parameterized output unify the PDE operators of BSM, CEV, and Heston into a single parameterized form. A single deep fully-connected network with 9-dimensional input simultaneously covers the solution mappings of all three models without training separate networks. The soft mask $\tanh(\xi/0.05)^2$ enables continuous interpolation between models: when $\xi\to 0$, it degrades to BSM/CEV; when $\xi$ is large, it activates the stochastic volatility terms of Heston. The additive parameterization uses the Black-Scholes analytical solution as a baseline, with the network learning absolute corrections, resolving the degeneracy of multiplicative parameterization in deep out-of-the-money regions. The Improved Constrained PINN (ICPINN) hard constraint technique is applied to exactly satisfy the terminal condition via output transformation.

In the LLM component, the DeepSeek API serves as an intelligent routing layer, achieving reliable parameter extraction and model selection through structured JSON output, supporting both Chinese and English natural language input.

Experimental results show that the unified PINN achieves MAE=0.067 on BSM (ATM relative error 0.31\%) and MAE=0.031 on Heston (ATM relative error 0.52\%), matching or outperforming independently trained PINN baselines. Notably, the independently trained Heston PINN completely fails to converge, while the unified framework successfully resolves this issue. The LLM routing layer achieves 100\% model selection accuracy on 20 test cases. Single inference takes approximately 2ms, far below the approximately 500ms of traditional finite difference methods.

\englishkeywords{Physics-Informed Neural Networks; Option Pricing; Black-Scholes-Merton Model; Heston Model; Large Language Models; Unified Framework}
